\subsection{Selected results}
\begin{center}\small
\begin{tabular}{lrrrrr}\toprule
Domain and variant & Variables & Degree & Rows & Terms & Height\\\midrule
$\Z$, dimension &278&4&188&12884&12227\\
$\Z$, balanced &279&4&189&5483&7200\\
$\Z$, sparse &279&4&189&5202&12662\\
$\Z$, fewer terms &284&4&194&4270&10864\\
$\Z$, sparsest listed &306&4&216&3870&14512\\\midrule
$\N$, balanced &229&4&179&5607&65892965\\
$\N$, smaller lower height &228&4&178&9522&560804902\\
$\N$, smaller sparse &228&4&178&7512&2445412949744\\\bottomrule
\end{tabular}
\end{center}
Here $\N=\{0,1,2,\ldots\}$ and $\Z$ denotes unrestricted integers.
The 279-variable balanced equation improves the incoming 282-variable
comparison on every reported varying measure: its parent has 192 rows,
5491 terms and height 128446. The 278-variable equation goes one variable
further at a support cost. The other integer comparisons are to the
matched 287-variable/197-row/4313-term and
309-variable/219-row/3922-term incoming equations, both of height 128446.
These are direct proof-calculus reductions, not absolute variable-count
records across all coefficient encodings.

\begin{center}
\begin{tabular}{lrrr}\toprule
Domain and branch & Variables & Exact degree & Binary operations\\\midrule
Strictly positive, $D=XY$, squared weight &10&626&77\\
Strictly positive, $D=XY$, product weight &10&464&78\\
Strictly positive, $D=q^4$, product weight &10&460&79\\
Strictly positive, partial normalization &10&384&80\\\midrule
$\Z$, auxiliary root and odd guard &14&86&112\\
$\Z$, strong root and odd guard &14&94&110\\
$\Z$, exact recompiled predecessor &13&156&125\\\bottomrule
\end{tabular}
\end{center}
The new 14-variable degree-86 equation improves both degree and evaluator
cost over the preceding degree-88/119-operation form. Degree 94 is a
lower-cost companion, not a degree improvement. A degree-88/115-operation
exact recompilation is also retained as a check. The 13-variable entry
preserves the entire earlier polynomial and only shortens its evaluator.
No smaller fixed-port payload or fully expanded coefficient height is
claimed for these circuit branches.
